\section{Nonlinearity and feature learning}
\label{sec:nonlinear-learning-proof}

We prove \Cref{thm:nonlinear-learning} using the real bounds already
established in \Cref{cp:fit}. All auxiliary notation below is local;
constants and $O(\cdot)$ bounds may depend on the entire fixed problem,
but not on width. The key simplification is to test feature increments
against their initial adjoints, rather than follow every feature acceleration.

\begin{lemma}[Nonlinear initial velocity]
\label{fl:input-lemma}
Every $Q^{(\ell)}$, $\ell\ge1$, is positive definite. The limiting initial
prediction velocity is nonaffine on the sphere, witnessed by an affine-annihilating
combination of four deterministic inputs chosen before initialization.
\end{lemma}
\begin{proof}
Write $v=x/\sqrt d$. Gaussian covariance recursion gives
$Q^{(\ell)}_{ab}=F_\ell(v_a^\top v_b)$, $F_0(s)=s$, with
\begin{equation}
 F_\ell(s)=\sum_{k\ge0}c_{\ell,k}s^k,\qquad
 c_{\ell,k}\ge0,\quad \sum_kc_{\ell,k}<\infty,
 \label{fl:positive-series}
\end{equation}
and unbounded positive coefficient support. Here is a direct justification.
The orthonormal Gaussian Hermite expansion of
$\phi_\ell(\sqrt qG)$ has coefficients $\alpha_k$ and covariance
$\sum_k\alpha_k^2r^k$ for Gaussian correlation $r$. The covariance identity
follows by comparing coefficients in
$\E[e^{zG-z^2/2}e^{wG'-w^2/2}]=e^{rzw}$.
Completeness follows because a function orthogonal to all polynomials has
entire transform $\E[h(G)e^{zG}]$ identically zero, hence zero Fourier
transform. Linear growth gives Gaussian $L^2$ integrability. Finite Hermite
support would make the activation a polynomial, hence affine by its bounded
derivative. Composing the covariance series with
$F_{\ell-1}(s)/F_{\ell-1}(1)$ proves \eqref{fl:positive-series}:
Tonelli at $s=1$ gives summability, and a positive inner coefficient of
degree $j\ge1$ and outer degree $k$ produces degree $jk$. These arguments
allow nonzero means and justify uniform absolute convergence on $[-1,1]$.
The tensor Grams $[(v_a^\top v_b)^k]$ are positive semidefinite and converge
to $I_m$ by the nonparallel-input hypothesis. Thus every $Q^{(\ell)}\succ0$.

Let $P$ be the orthogonal projection onto affine functions in spherical
$L^2$, with normalized surface measure $\sigma$:
$Ph(v)=\int h\,d\sigma+d\,v^\top\int uh(u)\,d\sigma(u)$.
Projecting the tensor feature in both variables gives the positive kernel
\[
 R_k(v,u)=(v^\top u)^k-a_k-b_kv^\top u,
 \quad a_k=\int(z^\top u)^k\,d\sigma(z),\quad
 b_k=d\int(z^\top u)^{k+1}\,d\sigma(z).
\]
Here $a_k,b_k\to0$ by dominated convergence, since $d\ge2$, so
$[R_k(v_a,v_b)]\to I_m$. The limiting velocity
$g(v)=(2/m)\sum_a y_aF_L(v^\top v_a)$ is not affine: otherwise
\[
 0=\sum_a y_a[(I-P)g](v_a)
   =\frac2m\sum_kc_{L,k}\,y^\top[R_k(v_a,v_b)]y>0.
\]
Every term is nonnegative and some large-degree term is positive.
Uniform convergence justifies the projection and sums.

Some great circle witnesses this nonaffinity. Otherwise the antipodal
average of $g$ is constant, and the odd remainder extends homogeneously
to a function linear on every plane, hence globally linear. On such a
circle, interpolate three distinct points by an affine function and take
a fourth where it fails. The resulting dependence can be signed and normalized
so that, for unit vectors $u_i$,
\begin{equation}
 \sum_i|c_i|=1,\qquad \sum_i c_i=0,\qquad
 \sum_i c_i u_i=0,\qquad \sum_i c_i g(u_i)>0.
 \label{fl:four-witness}
\end{equation}
These inputs depend on the fixed problem only. The finite-query Gram
convergence proved in \Cref{cp:fit}, applied to this augmented list, and
$\dot f_n(0,x)=(2/mn)\sum_a y_a h_a^{(L)}(0)^\top h^{(L)}(0,x)$
give convergence of their initial velocities.
\end{proof}

\begin{lemma}[Initial backward fields and hidden forces]
\label{fl:initial-activity}
The initialized forward fields and first backward derivatives, including
the vectors before derivative gates, have deterministic joint empirical
row-law limits with second moments. This includes the full first-layer
weight and acceleration rows. Every hidden block's squared initial
acceleration in inverse-mobility norm converges to a positive constant.
\end{lemma}
\begin{proof}
All unmarked fields in this proof are at zero. Put $k=2/m$ and define
the pre-gated vectors $P$ and their gated versions $B$ by
\[
 P_a^{(L)}=\sum_b y_bh_b^{(L)},\qquad
 B_a^{(\ell)}=\phi_\ell'(z_a^{(\ell)})\od P_a^{(\ell)},\qquad
 P_a^{(\ell-1)}=W^{(\ell)\top}B_a^{(\ell)}.
\]
Thus $\dot\delta_a^{(\ell)}(0)=kB_a^{(\ell)}$. Differentiating
\eqref{cp:flow}, since all hidden velocities initially vanish, gives
\begin{equation}
 \ddot W^{(1)}(0)=k^2\sum_a y_aB_a^{(1)}v_a^\top,\qquad
 \ddot W^{(\ell)}(0)=\frac{k^2}{n}\sum_a y_aB_a^{(\ell)}h_a^{(\ell-1)\top}.
 \label{fl:initial-forces}
\end{equation}

Only a forward pass and a backward pass are needed for their joint limits.
For one hidden Gaussian matrix $W$, after revealing $WH=Z$, its conditional
law is $Z(H^\top H)^{-1}H^\top+\widetilde W(I-\Pi_H)$, where
$\Pi_H=H(H^\top H)^{-1}H^\top$ and $\widetilde W$ is fresh with
entries $N(0,1/n)$. This is orthogonal Gaussian conditioning on the
linear constraint. For a reverse query $U$, known from the previous answers,
\begin{equation}
 W^\top U\ \overset d=\ H Q_n^{-1}C_n+(I-\Pi_H)\Xi D_n^{1/2},
 \quad Q_n=H^\top H/n,\ C_n=Z^\top U/n,\ D_n=U^\top U/n,
 \label{fl:reverse-conditioning}
\end{equation}
with independent standard Gaussian rows in $\Xi$. The removed projection
has conditional squared RMS
$\operatorname{rank}(H)\operatorname{tr}(D_n)/n$, hence vanishes in
probability. Adaptive queries cause no extra dependence: conditional on the
global transcript, the next query is fixed; revealing one Gaussian action
leaves the other matrices' residuals independent. In forward-then-backward
order no matrix is queried again after its reverse answer.

For precision, the induction uses empirical quadratic Wasserstein convergence.
Appending independent Gaussian rows preserves a deterministic such limit:
conditional variances for bounded tests are $O(1/n)$, and the appended
second moment obeys the law of large numbers. Continuous coordinate maps
of at most linear growth preserve this convergence, by an $L^2$ coupling
and uniform integrability. These include $\phi(z)$ and $\phi'(z)u$.
Same-layer products consequently converge; inverse Grams and positive
covariance square roots converge when their limiting Grams are, respectively,
positive definite or positive semidefinite. The vanishing projected term
in \eqref{fl:reverse-conditioning} does not change the row-law limit.
Starting with the full iid first-layer rows, these rules prove the forward
limits and then the backward limits, including all pre-gated fields.
The forward Grams being inverted are positive by \Cref{fl:input-lemma};
$Q^{(0)}$ is never inverted. The first-row acceleration is a linear
combination of the $B$ fields, so its joint limit is included.

Let $D^{(\ell)}$ be the limiting Gram of $B^{(\ell)}$.
At the top it is positive definite: a null vector $c$ would give
\[
 \Big(\sum_b y_b\phi_L(z_b)\Big)
 \Big(\sum_a c_a\phi_L'(z_a)\Big)=0\qquad(z\in\R^m).
\]
Indeed $L\ge2$ gives full Gaussian support. The first factor is not
identically zero since $y^\top Q^{(L)}y>0$; analyticity forces the
second to vanish everywhere. Varying each coordinate, the nonconstant
derivative gives $c=0$. Downward, the fresh reverse innovation in
\eqref{fl:reverse-conditioning} is independent of the lower-layer fields,
with covariance $D^{(\ell+1)}$, so
\[
 c^\top D^{(\ell)}c\ge\lambda_{\min}(D^{(\ell+1)})
       \sum_a c_a^2\E[\phi_\ell'(Z_a^{(\ell)})^2]>0.
\]
Each Gaussian marginal is nondegenerate and a nonzero analytic derivative
has a zero set of measure zero. Finally, expanding \eqref{fl:initial-forces},
the layer-$\ell$ squared acceleration limit is
$k^4y^\top(Q^{(\ell-1)}\circ D^{(\ell)})y>0$:
for $Q\succeq0$, $D\succ0$,
$Q\circ D\succeq\lambda_{\min}(D)\operatorname{diag}(Q)$ by decomposing
$Q$ into rank-one terms, and every
relevant $Q$ has positive diagonal. This also covers singular input Grams.
\end{proof}

\begin{lemma}[Adjoint test of the actual motion]
\label{fl:strong-time}
Every hidden layer satisfies \eqref{fl:layer-motion}. On the training set,
\begin{equation}
 y^\top(f_n(t)-f_{\mathrm{NTK},n}(t))
  =\frac m3t^3\left[\frac{\|\ddot W^{(1)}(0)\|_F^2}{n}
                 +\sum_{\ell=2}^L\|\ddot W^{(\ell)}(0)\|_F^2\right]+o_*(t^3).
 \label{fl:cubic-identity}
\end{equation}
Here $o_*(t^j)$ means: for every $\varepsilon>0$, some deterministic
$T_\varepsilon>0$ makes the probability that the normalized remainder
exceeds $\varepsilon$ anywhere on $0<t\le T_\varepsilon$ tend to zero.
\end{lemma}
\begin{proof}
Use $\|u\|_n=\|u\|_2/\sqrt n$ and the inverse-mobility norm
$\|U\|_{\rm mob}^2=\|U^{(1)}\|_F^2/n+\sum_{\ell\ge2}\|U^{(\ell)}\|_F^2$
on hidden matrices. On the event of \Cref{cp:fit}, all relevant operators,
feature RMS and residuals have deterministic bounds. The readout equation
therefore gives $\|w(t)\|_n=O(t)$; backpropagation gives
$\|\delta_a^{(\ell)}(t)\|_n=O(t)$; integration of the hidden updates and
forward subtraction give, uniformly on the sphere for $0\le t\le1$,
\begin{equation}
 \|W(t)-W(0)\|_{\rm mob}=O(t^2),\qquad
 \|\Delta z^{(\ell)}\|_n+\|\Delta h^{(\ell)}\|_n=O(t^2),\qquad
 f_n(t,x)=t\dot f_n(0,x)+O(t^2).
 \label{fl:real-bounds}
\end{equation}
Here $\Delta$ means the increment from zero. For the last assertion,
the readout equation gives $\|\dot w(t)-\dot w(0)\|_n=O(t)$; integrate
and use $f=w^\top h/n$. No new bootstrap is needed.

The remainder control uses one elementary estimate. For bounded Lipschitz
$g$ and an initial vector $u$,
\begin{equation}
 \|[g(z)-g(z_0)]\od u\|_n^2
 \le\operatorname{Lip}(g)^2D^2\|z-z_0\|_n^2
       +4\|g\|_\infty^2\|u\mathbf1_{|u|>D}\|_n^2.
 \label{fl:multiplier}
\end{equation}
Split at $|u_i|=D$. The deterministic row-law limits with second moments
in \Cref{fl:initial-activity} give, for every $\eta>0$, a fixed cutoff
$D$ for which the probability that any needed initial squared tail exceeds
$\eta$ tends to zero. Choose $D$ first, then a small time in
\eqref{fl:real-bounds}. Thus gate differences times initial fields are
$o_*(1)$, uniformly also along the segment from $z_0$ to $z$.

For this proof define $k=2/m$, $P_a^{(L)}=\sum_b y_bh_b^{(L)}(0)$,
$B_a^{(\ell)}=\phi_\ell'(z_a^{(\ell)}(0))\od P_a^{(\ell)}$, and
$P_a^{(\ell-1)}=W^{(\ell)}(0)^\top B_a^{(\ell)}$.
Since $w(t)/t=\dot w(0)+O_{\|\cdot\|_n}(t)$, downward subtraction of
the backward recursion yields
\begin{equation}
 \delta_a^{(\ell)}(t)=ktB_a^{(\ell)}+o_{*,\|\cdot\|_n}(t),\qquad
 \Delta W=\tfrac12t^2\ddot W(0)+o_{*,\|\cdot\|_{\rm mob}}(t^2).
 \label{fl:weight-expansion}
\end{equation}
Indeed each backward difference is the next-layer error through bounded
operators, an $O(t^2)$ operator increment acting on $B$, and a gate
difference acting on $P$, controlled by \eqref{fl:multiplier}.
Insert this into the hidden updates with $r=-y+O(t)$ and
$h=h(0)+O_{\|\cdot\|_n}(t^2)$, then integrate. The normalized outer-product
identities $\|uv^\top\|_F/\sqrt n=\|u\|_n\|v\|_2$ and
$\|uh^\top\|_F/n=\|u\|_n\|h\|_n$ justify the matrix-norm remainder.

Instead of expanding each forward feature, test its increment against $P$.
Put $E_{1,n}=\|\ddot W^{(1)}(0)\|_F^2/n$ and
$E_{\ell,n}=\|\ddot W^{(\ell)}(0)\|_F^2$ for $\ell\ge2$.
The integral formula for an activation increment and \eqref{fl:multiplier}
give
\[
 \langle P_a^{(\ell)},\Delta h_a^{(\ell)}\rangle_n
  =\langle B_a^{(\ell)},\Delta z_a^{(\ell)}\rangle_n+o_*(t^2),
 \qquad \langle u,v\rangle_n=u^\top v/n.
\]
The error is bounded by the RMS gate difference times $P$, multiplied by
$\|\Delta z\|_n=O(t^2)$. In
$\Delta z=\Delta W h(0)+W(0)\Delta h+\Delta W\Delta h$,
the last term is $O(t^4)$ in RMS. Summing the first term's pairings with
$y_aB_a$ and using \eqref{fl:initial-forces} gives
$t^2E_{\ell,n}/(2k^2)+o_*(t^2)$; the second pairs the lower feature
increment with $P_a^{(\ell-1)}$. Telescoping to the first layer yields
\begin{equation}
 \sum_a y_a\langle P_a^{(\ell)},\Delta h_a^{(\ell)}\rangle_n
   =\frac{t^2}{2k^2}\sum_{j\le\ell}E_{j,n}+o_*(t^2).
 \label{fl:adjoint-increment}
\end{equation}
All energies have positive limits. Cauchy--Schwarz, with bounded
$\sum_a y_a^2\|P_a^{(\ell)}\|_n^2$, therefore proves the feature-motion
lower bound in every layer on one deterministic small interval.

The tangent kernel is the readout Gram plus the hidden terms
\[
 K_{ab}(t)=\langle h_a^{(L)},h_b^{(L)}\rangle_n
 +(v_a^\top v_b)\langle\delta_a^{(1)},\delta_b^{(1)}\rangle_n
 +\sum_{\ell=2}^L\langle\delta_a^{(\ell)},\delta_b^{(\ell)}\rangle_n
                 \langle h_a^{(\ell-1)},h_b^{(\ell-1)}\rangle_n.
\]
Equation \eqref{fl:real-bounds} gives $K(t)-K(0)=O(t^2)$.
Its readout part, contracted with $y$, is twice the top pairing in
\eqref{fl:adjoint-increment}, plus $O(t^4)$.
By \eqref{fl:weight-expansion} and the squared outer products in
\eqref{fl:initial-forces}, the hidden contribution has the same coefficient.
Consequently
\[
 y^\top[K(t)-K(0)]y
       =\frac{2t^2}{k^2}\sum_\ell E_{\ell,n}+o_*(t^2).
\]
Finally $\dot f_n=kK(t)(y-f_n)$, whereas the frozen flow uses $K(0)$.
Variation of constants gives their difference as
$k\int_0^t e^{-kK(0)(t-s)}[K(s)-K(0)](y-f_n(s))\,ds$.
Here $f_n(s)=O(s)$ and the exponential is $I+O(t-s)$ because $K(0)$
is bounded. Pair with $y$ and integrate the last display. The additional
terms are $O(t^4)$, giving $2/(3k)=m/3$ in \eqref{fl:cubic-identity}.
\end{proof}

\begin{lemma}[The nonlinear feature component moves]
\label{fl:nonlinear-feature}
Let $V=\operatorname{span}\{x_1,\ldots,x_m\}$, with normalized surface
measure $\sigma_V$ on its unit sphere. For sufficiently small fixed $t>0$,
with probability tending to one,
\begin{equation}
 \frac1n\sum_i\inf_{a\in V,\,b\in\R}
 \int_{S(V)}|h_i^{(1)}(t,\sqrt d\,v)-h_i^{(1)}(0,\sqrt d\,v)
                  -a^\top v-b|^2\,d\sigma_V(v)\ge c_*t^4.
 \label{fl:nonlinear-increment}
\end{equation}
\end{lemma}
\begin{proof}
Since $\dim V\ge2$, for nonzero $g,r\in V$ the function
$\phi_1'(g^\top v)(r^\top v)$ is not affine on $S(V)$.
Otherwise its values on $r^\perp$ and antipodes force the affine function
to be $\lambda r^\top v$. Off the equator, division and then continuity
make $\phi_1'$ constant on $[-\|g\|,\|g\|]$, contradicting analyticity
and nonaffineness. Project the initial first-layer row onto $V$ and call it
$g_i$; call its acceleration row $r_i\in V$. By \Cref{fl:initial-activity},
their joint empirical law converges with second moments; $g$ is nondegenerate
Gaussian and $\E\|r\|^2>0$. The squared distance of
$\phi_1'(g^\top v)(r^\top v)$ from affine functions in $L^2(\sigma_V)$
is continuous in $(g,r)$, bounded by $\|\phi_1'\|_\infty^2\|r\|^2$,
and positive when both are nonzero. Its empirical mean therefore tends to
a positive constant.

The first-layer expansion in \eqref{fl:weight-expansion} and bounded
$\phi_1'$ reduce the actual increment to the fixed row direction
$t^2r_i/2$. On $\|r_i\|\le D$, its Taylor remainder divided by $t^2$
is at most $\|\phi_1''\|_\infty t^2D^2/8$; on the complement it is
at most $\|\phi_1'\|_\infty\|r_i\|$. The deterministic second-moment
limit, choosing $D$ and then time, gives $o_*(t^2)$ in neuron-index-times-sphere
$L^2$. Orthogonal projection off affine functions is contractive, so the
positive limiting projected acceleration proves the claim.
\end{proof}

\begin{proof}[Proof of \Cref{thm:nonlinear-learning}]
The four-point combination \eqref{fl:four-witness}, initial velocity
convergence and \eqref{fl:real-bounds} give a positive multiple of $t$
for the dense predictor on a width-independent small interval.
It annihilates every affine predictor and has coefficient absolute sum one,
so it bounds their uniform distance below. The positive energy limit in
\eqref{fl:cubic-identity}, divided by $\|y\|_1$, similarly gives a
training-input gap of order $t^3$ from the frozen kernel.
The two preceding lemmas supply all internal feature conclusions.

For each fixed positive time, the vanishing absolute compression errors
in \Cref{cp:legendre,cp:harmonic,cp:panel} transfer both prediction gaps
by the triangle inequality; eventual estimates at every confidence imply
convergence in probability. Choose common smaller $c_*,t_*>0$.
Legendre and Harmonic already cover the witnesses. Taylor appends them
before initialization, without labels. Its bound replaces $p$ by at most
$p+4$, multiplying the quadratic panel factor by at most $9$ and the
linear data factor by at most $3$, since $m+p\ge2$.
No endpoint gap or identification of compressed coordinates with dense
neurons is used or asserted.
\end{proof}
